\begin{table*}
  \caption{Tempo de execução (média, desvio-padrão e IC de 95\%) na maior componente ($n=1893$, $m=13835$) e complexidade teórica. IC pela Normal quando $n\ge30$ e pela t-Student quando $n<30$.}
  \label{tab:tempos}
  \small
  \begin{tabular}{llrrrl}
    \toprule
    Algoritmo & Complexidade & Rep. & Média & DP & IC 95\% \\
    \midrule
    BFS & $O(V + E)$ & 30 & 1,09 ms & 0,17 ms & [1,03 ms; 1,15 ms] \\
    DFS & $O(V + E)$ & 30 & 1,77 ms & 0,38 ms & [1,64 ms; 1,91 ms] \\
    Eulerianidade & $O(V + E)$ & 30 & 1,37 ms & 0,22 ms & [1,29 ms; 1,45 ms] \\
    Dijkstra & $O((V + E) \log V)$ & 30 & 4,32 ms & 0,54 ms & [4,13 ms; 4,51 ms] \\
    Bellman-Ford & $O(V\cdot E)$ & 30 & 14,00 ms & 1,22 ms & [13,56 ms; 14,43 ms] \\
    Tarjan (SCC) & $O(V + E)$ & 30 & 2,82 ms & 0,27 ms & [2,73 ms; 2,92 ms] \\
    Tarjan (pontes) & $O(V + E)$ & 30 & 3,80 ms & 0,39 ms & [3,66 ms; 3,94 ms] \\
    Kruskal & $O(E \log E)$ & 30 & 7,88 ms & 0,79 ms & [7,60 ms; 8,16 ms] \\
    Prim & $O(E \log V)$ & 30 & 13,27 ms & 0,95 ms & [12,94 ms; 13,61 ms] \\
    Floyd-Warshall & $O(V^3)$ & 5 & 23,62 s & 1,25 s & [22,07 s; 25,18 s] \\
    Floyd-Warshall ($n=150$, python) & $O(V^3)$ & 10 & 206,25 ms & 34,04 ms & [181,91 ms; 230,60 ms] \\
    Floyd-Warshall ($n=150$, numpy) & $O(V^3)$ & 10 & 4,48 ms & 0,17 ms & [4,36 ms; 4,61 ms] \\
    \bottomrule
  \end{tabular}
\end{table*}
